


\documentclass[11pt]{article}
\usepackage[utf8]{inputenc}
\usepackage[english]{babel}
\usepackage[round]{natbib}
\setcitestyle{authoryear,open={(},close={)}} 


\usepackage{mathpazo, amsmath, amssymb, amsfonts, amsthm, microtype, tikz, titling, graphicx, booktabs, caption, subcaption, bm, xfrac, appendix, setspace, comment, float, nth, fnpct, makecell, multirow, tabularx, threeparttable, titling, array, lipsum, mathtools, tocloft}


%\bibliographystyle{abbrvnat}
\usepackage{enumitem}
\usepackage{graphics}
\usepackage{url}
\linespread{1.5}
\usepackage[margin=0.8in]{geometry}
\usepackage{hyperref}
\usepackage{comment}

\hypersetup{
    colorlinks,
    linkcolor={blue!80!black},
    citecolor={blue!90!black},
    urlcolor={blue!90!black}, 
   % urlbordercolor=red,
   % linkbordercolor=red,
   % citebordercolor=red,
   % runbordercolor=red,
   % citebordercolor=red,
    %filebordercolor=red
}


\usepackage{xcolor}
\usepackage{graphicx}
\usepackage{url}
\usepackage[round]{natbib}
\usepackage{amsmath,amsthm} 
\usepackage{engord}
\usepackage{float}
\usepackage{subfig}
\usepackage{pdflscape}
\usepackage{booktabs}
\usepackage{pgfplots}
\usepackage{titleps}
\usepackage{setspace}
\usepackage{dcolumn}
%\usepackage{subcaption}

%references
\usepackage{cleveref}

%%%%%%%
\usepackage{blindtext}
\usepackage{tabularx}
\usepackage{bbold}
%%%%%%%
\usepackage{multirow}
\usepackage{xr}
\usepackage{setspace}
\usepackage{sectsty}
\usepackage[labelsep=period]{caption} %% This switches %"Table 1: itle" to "Table 1. Title"
\usepackage{rotating}
\usepackage{graphics}
\usepackage{fancyhdr, ifthen}
\usepackage{blindtext}
\usepackage{pdflscape}

	
\usepackage{lineno}
%\linenumbers


\usepackage{amssymb}
\newtheorem{prediction}{Prediction}

\usepackage{appendix}
\newcommand{\Appendixautorefname}{appendix}

%%%%%%%%%%%%%%%%%%%%%%%%%%%%%%%%%%%%%%%%%

\newcommand{\footremember}[2]{%
    \footnote{#2}
    \newcounter{#1}
    \setcounter{#1}{\value{footnote}}%
}

\newcommand{\footrecall}[1]{%
    \footnotemark[\value{#1}]%
} 

% Adjust vertical space between main text and footnotes
\skip\footins=2\baselineskip plus 4pt minus 2pt

%\setlength{\footnotesep}{\baselineskip}

\usepackage{multicol}

\usepackage{threeparttable, booktabs, adjustbox}

\usepackage{subfig} % Required for subfigures

\usepackage[bottom]{footmisc} % Ensure footnotes are always at the bottom

%\flushbottom
%\raggedbottom
%\interfootnotelinepenalty=10000
\interfootnotelinepenalty=10000 % Set a high penalty




%%%%%%%%%%%%%%%%%% TITLE %%%%%%%    %%%%%%%%

\title{\centering The Political Legacy of Displacement: \\ Evidence from the Spanish Exile in France}
\author{
Cristina Aranzana\footnote{University of Barcelona. Email: \href{mailto:c.aranzana@ub.edu}{c.aranzana@ub.edu}} \and
Luc Paluskiewicz\footnote{Paris School of Economics. Email: \href{mailto:luc.paluskiewicz@psemail.eu}{luc.paluskiewicz@psemail.eu}}
}

\date{Extended Abstract \\ \today. } 

%%%%%%%%%%%%%%%%%% TEXT %%%%%%%%%%%%%%%

\begin{document}

\maketitle 

\\ 
\\
\vspace{0.1in}
\textbf{Session:} Guerres, reconstructions et industries. \\ 
\vspace{0.1in}

In early 1939, the defeat of Republican forces in Catalonia during the Spanish Civil War precipitated a mass exodus known as ``La Retirada" (``The Retreat”). Between late January and mid-February 1939, close to half a million Spanish Republicans fled across the Pyrenees into France. In a context of political tensions, French authorities, unprepared for the magnitude of the flow, treated the influx as a security concern and responded by interning the displaced population. Across continental France, more than 900 centres and concentration camps were created to house Spanish refugees, including both soldiers and civilians. Due to the magnitude of the refugee flows, the locations of these camps were chosen under extreme time pressure and mainly based on logistical availability. \\

This paper aims to answer the following research question: How can political refugees influence the political and military mobilization of native populations? We study how the presence of Spanish political refugees enhanced short-term mobilization into Resistance movements during the Second World War. In order to do so, we exploit the quasi-random placement of these camps to study the political consequences of forced displacement. We compare the involvement of the local population in the Resistance in towns that hosted refugees to municipalities that never hosted Spanish refugees. We use a large set of different data sources, including newly digitized archival records on Spanish refugees, as well as data on French Resistance members and electoral data at the municipality level. \\

Our preliminary results suggest that exposure to the Spanish refugee population has a positive effect on enrolment in French Resistance organizations. We find that in municipalities close to refugee camps, the number of members of the Resistance is higher compared to municipalities that were not exposed to Spanish refugees. We argue that contact with foreign refugees, who were highly politically involved before leaving Spain and fleeing an authoritarian regime, increased participation in the Resistance. We show that the effects we measured are highly localized and that municipalities hosting camps or located within a few kilometres exhibit a larger positive effect. We also observe stronger effects among younger cohorts, consistent with the idea that younger populations are more influenced by their proximity to foreign political refugees. \\

Additionally, we study the settlement of this population and its long-term effects on the transmission of political values and social mobilization between the end of WWII and today. We further investigate the mechanisms of political backlash and idea diffusion, and assess the persistence of these dynamics by analysing the settlement of Spanish communities and their collective mobilization. \



%%%%%%%%%%%%%%%%%%%%%%%%%%%%%%%%%%%%%%%%%%%%%%%%%
\clearpage
\begin{singlespace}
%\bibliographystyle{plainnat}
\bibliographystyle{chicago}
\bibliographystyle{aer}
\bibliography{biblio}
\end{singlespace}
%%%%%%%%%%%%%%%%%%%%%%%%%%%%%%%%%%%%%%%%%%%%%%%%%



\end{document}

